% Options for packages loaded elsewhere
\PassOptionsToPackage{unicode}{hyperref}
\PassOptionsToPackage{hyphens}{url}
\PassOptionsToPackage{dvipsnames,svgnames,x11names}{xcolor}
\PassOptionsToPackage{space}{xeCJK}
%
\documentclass[
  10pt,
  a4paper,
]{article}
\usepackage{amsmath,amssymb}
\usepackage{setspace}
\usepackage{iftex}
\ifPDFTeX
  \usepackage[T1]{fontenc}
  \usepackage[utf8]{inputenc}
  \usepackage{textcomp} % provide euro and other symbols
\else % if luatex or xetex
  \usepackage{unicode-math} % this also loads fontspec
  \defaultfontfeatures{Scale=MatchLowercase}
  \defaultfontfeatures[\rmfamily]{Ligatures=TeX,Scale=1}
\fi
\usepackage{lmodern}
\ifPDFTeX\else
  % xetex/luatex font selection
    \setmainfont[]{Times New Roman}
    \setsansfont[]{Arial}
    \setmonofont[]{Menlo}
  \ifXeTeX
    \usepackage{xeCJK}
    \setCJKmainfont[]{Songti SC}
          \fi
  \ifLuaTeX
    \usepackage[]{luatexja-fontspec}
    \setmainjfont[]{Songti SC}
  \fi
\fi
% Use upquote if available, for straight quotes in verbatim environments
\IfFileExists{upquote.sty}{\usepackage{upquote}}{}
\IfFileExists{microtype.sty}{% use microtype if available
  \usepackage[]{microtype}
  \UseMicrotypeSet[protrusion]{basicmath} % disable protrusion for tt fonts
}{}
\makeatletter
\@ifundefined{KOMAClassName}{% if non-KOMA class
  \IfFileExists{parskip.sty}{%
    \usepackage{parskip}
  }{% else
    \setlength{\parindent}{0pt}
    \setlength{\parskip}{6pt plus 2pt minus 1pt}}
}{% if KOMA class
  \KOMAoptions{parskip=half}}
\makeatother
\usepackage{xcolor}
\usepackage[top=20mm,bottom=19mm,left=20mm,right=20mm]{geometry}
\usepackage{longtable,booktabs,array}
\usepackage{calc} % for calculating minipage widths
% Correct order of tables after \paragraph or \subparagraph
\usepackage{etoolbox}
\makeatletter
\patchcmd\longtable{\par}{\if@noskipsec\mbox{}\fi\par}{}{}
\makeatother
% Allow footnotes in longtable head/foot
\IfFileExists{footnotehyper.sty}{\usepackage{footnotehyper}}{\usepackage{footnote}}
\makesavenoteenv{longtable}
\setlength{\emergencystretch}{3em} % prevent overfull lines
\providecommand{\tightlist}{%
  \setlength{\itemsep}{0pt}\setlength{\parskip}{0pt}}
\setcounter{secnumdepth}{-\maxdimen} % remove section numbering
\usepackage{xeCJK}
\setCJKmainfont{Songti SC}
\setCJKsansfont{Heiti SC}
\setCJKmonofont{Heiti SC}
\usepackage{graphicx,multicol,placeins,tocloft}
\usepackage{titlesec,needspace,seqsplit,fancyhdr,caption,fvextra,etoolbox}
\definecolor{ReportNavy}{HTML}{18384C}
\definecolor{ReportTeal}{HTML}{236A78}
\definecolor{ReportGray}{HTML}{657582}
\definecolor{ReportRule}{HTML}{D9E2E7}
\definecolor{ReportCode}{HTML}{F3F6F8}
\AtBeginDocument{\hypersetup{colorlinks=true,linkcolor=ReportNavy,urlcolor=ReportTeal,bookmarksopen=true,bookmarksdepth=2,pdftitle={Lynx S10 地形感知自主导航与残差越障控制},pdfauthor={GOAI LAI / 狗來}}}
\pagestyle{fancy}
\fancyhf{}
\fancyhead[L]{\sffamily\fontsize{8}{10}\selectfont\color{ReportGray}GOAI LAI / Lynx S10}
\fancyhead[R]{\sffamily\fontsize{8}{10}\selectfont\color{ReportGray}TECHNICAL SYSTEM REPORT}
\fancyfoot[L]{\sffamily\fontsize{8}{10}\selectfont\color{ReportGray}2026-09-02}
\fancyfoot[R]{\sffamily\fontsize{8}{10}\selectfont\color{ReportGray}\thepage}
\renewcommand{\headrulewidth}{0.35pt}
\renewcommand{\footrulewidth}{0pt}
\setlength{\headheight}{15pt}
\setlength{\parindent}{0pt}
\setlength{\parskip}{5pt plus 1pt minus 1pt}
\setlength{\emergencystretch}{3em}
\widowpenalty=10000
\clubpenalty=10000
\displaywidowpenalty=10000
\raggedbottom
\titleformat{\section}{\sffamily\Large\bfseries\color{ReportNavy}}{}{0pt}{}
\titleformat{\subsection}{\sffamily\normalsize\bfseries\color{ReportNavy}}{}{0pt}{}
\titlespacing*{\section}{0pt}{18pt}{8pt}
\titlespacing*{\subsection}{0pt}{12pt}{5pt}
\pretocmd{\section}{\Needspace{6\baselineskip}}{}{}
\pretocmd{\subsection}{\Needspace{5\baselineskip}}{}{}
\captionsetup{font=small,labelfont=bf,justification=centering,skip=5pt}
\renewcommand{\figurename}{图}
\renewcommand{\topfraction}{0.92}
\renewcommand{\textfraction}{0.06}
\renewcommand{\floatpagefraction}{0.60}
\renewcommand{\arraystretch}{1.22}
\setlength{\tabcolsep}{4pt}
\AtBeginEnvironment{longtable}{\fontsize{8.5}{12.6}\selectfont}
\AtEndEnvironment{longtable}{\normalsize}
\DefineVerbatimEnvironment{ReportCode}{Verbatim}{fontsize=\footnotesize,breaklines=true,breakanywhere=true,frame=single,rulecolor=\color{ReportRule},framesep=5pt,baselinestretch=1.12}
\setcounter{tocdepth}{2}
\setlength{\cftbeforesecskip}{5pt}
\setlength{\cftbeforesubsecskip}{1pt}
\setlength{\cftsubsecindent}{1em}
\renewcommand{\cftsecfont}{\sffamily\small\bfseries}
\renewcommand{\cftsubsecfont}{\small}
\renewcommand{\cftsecpagefont}{\small}
\renewcommand{\cftsubsecpagefont}{\small}
\renewcommand{\contentsname}{目录}
\newcommand{\reporttitle}{\begingroup\sffamily\color{ReportNavy}\fontsize{24}{32}\selectfont\bfseries Lynx S10\par 地形感知自主导航与残差越障控制\par\endgroup\vspace{7pt}{\sffamily\color{ReportTeal}\large GOAI LAI / 狗來 · Technical System Report}\par\vspace{10pt}\textcolor{ReportRule}{\rule{\linewidth}{0.8pt}}\vspace{5pt}}
\ifLuaTeX
  \usepackage{selnolig}  % disable illegal ligatures
\fi
\usepackage{bookmark}
\IfFileExists{xurl.sty}{\usepackage{xurl}}{} % add URL line breaks if available
\urlstyle{same}
\hypersetup{
  colorlinks=true,
  linkcolor={Maroon},
  filecolor={Maroon},
  citecolor={Blue},
  urlcolor={Blue},
  pdfcreator={LaTeX via pandoc}}

\author{}
\date{}

\begin{document}

\setstretch{1.18}
\reporttitle

\textbf{系统基线：}
{\ttfamily\small\seqsplit{goai26-s10-racing@3660b81e8244dfa238633c161673596fe91650f6}}\\
\textbf{上游基线：}
{\ttfamily\small\seqsplit{goai\_embodied\_future\_material@13dd084be6cb5e2514098bc87e586d00dfe580b2}}\\
\textbf{结果基线：} 2026-08-20，Mac ARM64，seed 8，WP0→WP32 连续仿真自测
392.257 s\\
\textbf{文档范围：} 已交付系统的实现、Gate 16
策略训练与部署、实验结果及复现边界。\\
\textbf{修订日期：} 2026-09-02

\FloatBarrier

\section{摘要}\label{ux6458ux8981}

本项目面向 GOAI 2026 Track 4 Challenge 2，在官方 Lynx S10 轮足模型与
MuJoCo 赛道上实现分层自主导航。系统以在线几何 ray casting 构造 8×64
LiDAR range image 和 13×9 yaw-aligned height
map，使用严格顺序航点管理、前视路径跟踪、候选航向局部规划及地形状态机生成速度指令。普通赛段采用官方
57D 本体感知 locomotion policy；WP15→WP16 的 0.3769 m 高台阶由冻结的
174D base actor 与 heightmap-gated residual actor 接管。规则型 strategy
router 约束入口状态，并在四轮几何净空、接触及姿态检查后，通过 SDK
内单一关节控制权仲裁切回官方策略。

部署 residual 的 checkpoint 为 {\ttfamily\small\seqsplit{speed\_core}}
第 120 次 PPO 迭代，actor 为
{\ttfamily\small\seqsplit{174--512--512--256--128--16}} ELU
MLP。其训练采用固定官方台阶、入口状态随机化、冻结
base、参考策略约束和速度导向奖励。连续系统自测按序完成全部 33
个航点，仿真计时 392.257
s。该结果依赖仿真真值定位、几何高度采样和轮子接触状态；当前没有真机完成率或独立场景泛化的充分证据。

\clearpage
{\fontsize{9}{12}\selectfont
\begin{multicols}{2}
\tableofcontents
\end{multicols}}
\clearpage

\section{1.
任务定义与系统范围}\label{ux4efbux52a1ux5b9aux4e49ux4e0eux7cfbux7edfux8303ux56f4}

给定顺序航点集合 \(W=(w_0,\ldots,w_{32})\)，系统生成机身速度命令
\(c_t=(v_x,v_y,\omega_z)\)，并通过运动策略产生 16
个执行器目标。目标是在满足顺序到点和运行约束的条件下减少完整路线用时。

官方 evaluator 使用机器人参考位置与航点的水平距离：

\[
d_{xy}(p_t,w_i)=\sqrt{(x_t-x_i)^2+(y_t-y_i)^2},\qquad r_{\mathrm{eval}}=0.20\ \mathrm{m}.
\]

Follower 与 router 各自维护同一航点序列的 cursor，均使用
\(r_{\mathrm{internal}}=0.18\ \mathrm{m}\)。后续航点不会补偿前序漏点。该阈值与碰撞模型、轮子半径、局部规划的机身走廊相互独立。

指定仿真路线有 33 个航点，航点折线水平长度 224.21 m，累计上升 6.70
m。路线坐标、部分平台几何、Gate 16
台沿及若干困难路段配置在运行前已知。系统属于\textbf{已知任务路线上的闭环导航与专用越障集成}，当前没有在线
SLAM、全局最优路径求解或未知场景端到端策略。

官方资源提供机器人模型、物理环境、基础 locomotion actor、SDK
和计时器；团队实现感知扩展、导航、越障策略及接口、控制权协调、恢复逻辑和评测工具。{[}S1--S3{]}

\FloatBarrier

\section{2.
系统架构与运行实现}\label{ux7cfbux7edfux67b6ux6784ux4e0eux8fd0ux884cux5b9eux73b0}

\subsection{2.1 控制结构}\label{ux63a7ux5236ux7ed3ux6784}

\begin{figure}[tbp]
\centering
\includegraphics[width=\linewidth,height=0.68\textheight,keepaspectratio]{report_assets/figure-1.pdf}
\caption{系统数据流与关节控制权仲裁。箭头表示状态、命令或控制权消息。}
\end{figure}

Router
是带阈值、计时器、迟滞和状态锁存的有限状态机，未训练策略选择网络。Locomotion
层使用神经网络；任务层和局部规划层采用确定性算法与场地配置。

\subsection{2.2 实现模块}\label{ux5b9eux73b0ux6a21ux5757}

以下路径均相对于系统基线仓库；完整本地源码索引见 §12。

\Needspace{6\baselineskip}
{\small\sffamily\bfseries\color{ReportGray}表 1. 实现模块与职责\par}
\vspace{-4pt}

\begin{longtable}[]{@{}
  >{\raggedright\arraybackslash}p{(\columnwidth - 4\tabcolsep) * \real{0.1700}}
  >{\raggedright\arraybackslash}p{(\columnwidth - 4\tabcolsep) * \real{0.4300}}
  >{\raggedright\arraybackslash}p{(\columnwidth - 4\tabcolsep) * \real{0.4000}}@{}}
\toprule\noalign{}
\begin{minipage}[b]{\linewidth}\raggedright
模块
\end{minipage} & \begin{minipage}[b]{\linewidth}\raggedright
文件
\end{minipage} & \begin{minipage}[b]{\linewidth}\raggedright
主要职责
\end{minipage} \\
\midrule\noalign{}
\endhead
\bottomrule\noalign{}
\endlastfoot
仿真与感知 &
{\ttfamily\small\seqsplit{src/s10\_perception/s10\_perception/sim\_node.py}}、{\ttfamily\small\seqsplit{lidar.py}}、{\ttfamily\small\seqsplit{heightmap.py}}
& 扩展上游仿真，发布测距、高度、机身与轮子状态 \\
导航 &
{\ttfamily\small\seqsplit{src/s10\_auto\_nav/s10\_auto\_nav/follower\_node.py}}
& 航点、地形、局部规划和恢复的执行顺序 \\
跟踪与避障 &
{\ttfamily\small\seqsplit{pure\_pursuit.py}}、{\ttfamily\small\seqsplit{local\_planner.py}}、{\ttfamily\small\seqsplit{waypoints.py}}
& 前视目标、速度命令、机身走廊和顺序到点 \\
地形通过 &
{\ttfamily\small\seqsplit{terrain.py}}、{\ttfamily\small\seqsplit{step\_commit.py}}
& 地形分类、持续通过和有界尝试 \\
策略协调 &
{\ttfamily\small\seqsplit{strategy\_router\_node.py}}、{\ttfamily\small\seqsplit{strategy/router.py}}
& 状态装配、入口约束、接管、退出与异常处理 \\
越障推理 &
{\ttfamily\small\seqsplit{integration/gate16\_policy\_runner.hpp}} &
174D observation、双 ONNX 推理、动作解码 \\
执行器仲裁 &
{\ttfamily\small\seqsplit{integration/joint\_command\_owner.hpp}} & 在
SDK 输出边界决定有效关节命令 \\
速度接口 &
{\ttfamily\small\seqsplit{integration/ros\_cmd\_interface.hpp}} & 接入
{\ttfamily\small\seqsplit{/cmd\_vel}}、起立与超时处理 \\
数据记录 & {\ttfamily\small\seqsplit{segment\_recorder.py}} &
逐时刻轨迹、动作、控制权和失败记录 \\
\end{longtable}

\subsection{2.3 ROS
与执行器接口}\label{ros-ux4e0eux6267ux884cux5668ux63a5ux53e3}

\Needspace{6\baselineskip}
{\small\sffamily\bfseries\color{ReportGray}表 2. ROS 与执行器接口\par}
\vspace{-4pt}

\begin{longtable}[]{@{}
  >{\raggedright\arraybackslash}p{(\columnwidth - 6\tabcolsep) * \real{0.2900}}
  >{\raggedright\arraybackslash}p{(\columnwidth - 6\tabcolsep) * \real{0.2500}}
  >{\raggedleft\arraybackslash}p{(\columnwidth - 6\tabcolsep) * \real{0.1300}}
  >{\raggedright\arraybackslash}p{(\columnwidth - 6\tabcolsep) * \real{0.3300}}@{}}
\toprule\noalign{}
\begin{minipage}[b]{\linewidth}\raggedright
Topic
\end{minipage} & \begin{minipage}[b]{\linewidth}\raggedright
类型
\end{minipage} & \begin{minipage}[b]{\linewidth}\raggedleft
名义频率
\end{minipage} & \begin{minipage}[b]{\linewidth}\raggedright
用途
\end{minipage} \\
\midrule\noalign{}
\endhead
\bottomrule\noalign{}
\endlastfoot
{\ttfamily\small\seqsplit{/ground\_truth/odom}} &
{\ttfamily\small\seqsplit{nav\_msgs/Odometry}} & 50 Hz &
Follower、router 和 recorder 的机身状态 \\
{\ttfamily\small\seqsplit{/scan}} &
{\ttfamily\small\seqsplit{sensor\_msgs/LaserScan}} & 50 Hz &
近水平障碍测距 \\
{\ttfamily\small\seqsplit{/perception/lidar}} &
{\ttfamily\small\seqsplit{Float32MultiArray}} & 50 Hz & 8×64
测距及可视化 \\
{\ttfamily\small\seqsplit{/perception/heightmap}} &
{\ttfamily\small\seqsplit{Float32MultiArray}} & 50 Hz & 117
个局部地形值 \\
{\ttfamily\small\seqsplit{/perception/wheel\_state}} &
{\ttfamily\small\seqsplit{Float32MultiArray}} & 50 Hz &
四轮位置与接触状态 \\
{\ttfamily\small\seqsplit{/strategy/nav\_cmd\_vel}} &
{\ttfamily\small\seqsplit{geometry\_msgs/Twist}} & 50 Hz &
Follower→router 的候选速度命令 \\
{\ttfamily\small\seqsplit{/cmd\_vel}} &
{\ttfamily\small\seqsplit{geometry\_msgs/Twist}} & 50 Hz & Router→SDK
速度接口，供当前运动策略使用 \\
{\ttfamily\small\seqsplit{/strategy/joint\_owner}} &
{\ttfamily\small\seqsplit{std\_msgs/String}} & 事件／控制周期 &
请求控制权 \\
{\ttfamily\small\seqsplit{/joints/owner}} &
{\ttfamily\small\seqsplit{std\_msgs/String}} & 事件 & SDK
确认的实际控制权 \\
{\ttfamily\small\seqsplit{/JOINTS\_DATA}}、{\ttfamily\small\seqsplit{/JOINTS\_CMD}}
& 上游自定义消息 & SDK 控制频率 & 实测关节状态与最终命令矩阵 \\
\end{longtable}

启用 router 时，launch 将 follower 输出重映射到
{\ttfamily\small\seqsplit{/strategy/nav\_cmd\_vel}}，由 router 统一发布
{\ttfamily\small\seqsplit{/cmd\_vel}}。Gate 16 模型在 SDK
进程内执行，通过同一 arbiter 输出关节命令；另有 ROS climb-command
接口用于可选策略，但不是本次 Gate 16 执行路径。

\FloatBarrier

\section{3.
感知与状态表示}\label{ux611fux77e5ux4e0eux72b6ux6001ux8868ux793a}

\subsection{3.1 合成 LiDAR}\label{ux5408ux6210-lidar}

{\ttfamily\small\seqsplit{mj\_multiRay}}
对当前碰撞几何执行在线查询。射线模板预计算，运行时按机身姿态变换，复用距离与
geom-id 缓冲区。

\Needspace{6\baselineskip}
{\small\sffamily\bfseries\color{ReportGray}表 3. 合成 LiDAR 配置\par}
\vspace{-4pt}

\begin{longtable}[]{@{}
  >{\raggedright\arraybackslash}p{(\columnwidth - 2\tabcolsep) * \real{0.2800}}
  >{\raggedright\arraybackslash}p{(\columnwidth - 2\tabcolsep) * \real{0.7200}}@{}}
\toprule\noalign{}
\begin{minipage}[b]{\linewidth}\raggedright
参数
\end{minipage} & \begin{minipage}[b]{\linewidth}\raggedright
配置
\end{minipage} \\
\midrule\noalign{}
\endhead
\bottomrule\noalign{}
\endlastfoot
方位采样 & 64，360°，端点不重复 \\
仰角采样 & 8，−30° 至 +15° 均匀采样 \\
测距范围 & 0.05--12.0 m \\
机身系安装偏置 & {\ttfamily\small\seqsplit{{[}0.20, 0, 0.12{]} m}} \\
几何过滤 & 地形 group 0；排除机器人和非碰撞航点标记 \\
未命中 & 发布 12.0 m \\
{\ttfamily\small\seqsplit{/scan}} 来源 & 最接近水平的仰角行，实际为约
+2.143° \\
\end{longtable}

该传感器通过射线与场景的首次相交建模几何遮挡；同一帧使用同一仿真状态，未模拟测量噪声、材质回波、扫描运动畸变或传输延迟。

\subsection{3.2
航向对齐高度图}\label{ux822aux5411ux5bf9ux9f50ux9ad8ux5ea6ux56fe}

局部网格在机身前后方向取 13 点、左右方向取 9 点：

\[
x\in[-0.60,1.20],\quad y\in[-0.60,0.60],\quad \Delta x=\Delta y=0.15\ \mathrm{m}.
\]

网格随 yaw 对齐，忽略机身 roll/pitch；展平采用 x 外循环、y
内循环。高度值通过网格位置的向下射线直接查询场景，\textbf{不是从 8×64
LiDAR 重建得到}。实现可跳过高于
{\ttfamily\small\seqsplit{base\_z+1.0 m}} 的上方表面，最多 12
次，以减少多层结构干扰。

导航使用：

\[
h^{\mathrm{raw}}_{ij}=\operatorname{clip}(z^{\mathrm{terrain}}_{ij}-z^{\mathrm{base}},-1,1).
\]

Gate 16 使用训练约定：

\[
h^{\mathrm{policy}}_{ij}=\operatorname{clip}(-h^{\mathrm{raw}}_{ij}-0.5,-1,1),
\]

但 raw≤−1 的无效／空洞编码保留为 −1。实际裁剪到 −1
的深地面与无命中共享数值，当前消息没有逐格 validity mask。SDK
高度图缓存对超过 200 ms 的观测标记无效；无效图不能新启动 residual gate。

\subsection{3.3
地形状态与特权信息}\label{ux5730ux5f62ux72b6ux6001ux4e0eux7279ux6743ux4fe1ux606f}

高度特征先在轮子通行走廊内拟合平面，再计算相对该平面的正／负高度残差、平面坡度和起伏覆盖宽度。这样将连续坡面与局部台阶分开；分类阈值中的
rise 是平面残差，不等于障碍的绝对高度。

{\ttfamily\small\seqsplit{TerrainClassifier}}
融合这些特征、前向净空和姿态，输出
{\ttfamily\small\seqsplit{FLAT/RAMP/STAIRS/BLOCKED/HIGH\_BARRIER/DROP/UNSTABLE/UNKNOWN}}，通过
dwell 与 hysteresis 抑制标签抖动。只有
{\ttfamily\small\seqsplit{RAMP/STAIRS}} 可直接授权 step
commitment；高障碍先触发局部重规划，落差与未知状态限制前进。

当前部署依赖的 simulator-only 信息包括：真值
odometry、网格位置的直接地形查询，以及用于交接确认的四轮世界坐标和接触
flags。模型输入、router 输入与训练奖励使用的信息范围不同，应分别评估
sim-to-real 替代方案。

\FloatBarrier

\section{4.
导航与地形条件速度控制}\label{ux5bfcux822aux4e0eux5730ux5f62ux6761ux4ef6ux901fux5ea6ux63a7ux5236}

\subsection{4.1
顺序跟踪与前视控制}\label{ux987aux5e8fux8ddfux8e2aux4e0eux524dux89c6ux63a7ux5236}

前视距离由上一周期前向命令决定：

\[
L_t=1.5+0.7|v_{x,t-1}|.
\]

对选择后的目标 \(p^*\)，令
\(\Delta=p^*-p_t\)、\(e_\psi=\operatorname{wrap}(\operatorname{atan2}(\Delta_y,\Delta_x)-\psi_t)\)，实现计算：

\[
\omega_z=\operatorname{clip}(1.8e_\psi,-0.7,0.7),
\] \[
v_y=\operatorname{clip}\{0.9[-\sin\psi_t\Delta_x+\cos\psi_t\Delta_y],-0.4,0.4\}.
\]

当 \(|e_\psi|>30^\circ\)
时目标平移速度为零，先旋转对齐；否则前向命令随朝向误差线性衰减，并叠加到点制动。上述是变化率限制前的目标命令；实际输出的三个通道分别受
5.0 m/s²、2.0 m/s²、6.0 rad/s²
限制，因此切入旋转模式不等于平移命令瞬时归零。横向项是目标点在机身系的侧向偏移。该实现未采用标准自行车模型的曲率公式。

部分宽转角采用 corner
preview：世界系平移仍指向当前未计分航点，机身朝向提前插值到出弯方向。当前配置仅对
WP13、19、20、21、22、23 启用，前视触发距离 0.8 m、速度上限 0.7
m/s、倾角上限 10°；顺序到点条件不变。

\subsection{4.2
候选航向局部规划}\label{ux5019ux9009ux822aux5411ux5c40ux90e8ux89c4ux5212}

规划器围绕目标方位，在 ±90° 内采样 31 个航向。将 LiDAR
回波投影到候选方向，取 0.45 m 半宽走廊内最近的正向距离
\(C(\theta)\)，截断于 4 m。候选得分为：

\[
J(\theta)=1.6\frac{C(\theta)}{4.0}-\frac{|\operatorname{wrap}(\theta-\theta_g)|}{\pi/2}.
\]

选择最高得分，平分时优先较小目标偏差。选中净空不足 0.8 m
则判为阻塞，否则按净空缩放速度。若目标很近，超过目标距离加 0.25 m
机身前伸余量的障碍不否决该目标，避免航点后方柱子导致永远无法到点。

该方法是贪心候选航向搜索，未求解完整时域动力学轨迹，也不保证全局最优或连续碰撞安全。

\subsection{4.3
速度调度与场地先验}\label{ux901fux5ea6ux8c03ux5ea6ux4e0eux573aux5730ux5148ux9a8c}

\Needspace{6\baselineskip}
{\small\sffamily\bfseries\color{ReportGray}表 4. 导航速度调度与场地先验\par}
\vspace{-4pt}

\begin{longtable}[]{@{}
  >{\raggedright\arraybackslash}p{(\columnwidth - 2\tabcolsep) * \real{0.3400}}
  >{\raggedright\arraybackslash}p{(\columnwidth - 2\tabcolsep) * \real{0.6600}}@{}}
\toprule\noalign{}
\begin{minipage}[b]{\linewidth}\raggedright
条件／目标航点
\end{minipage} & \begin{minipage}[b]{\linewidth}\raggedright
命令上限或处理
\end{minipage} \\
\midrule\noalign{}
\endhead
\bottomrule\noalign{}
\endlastfoot
WP2、10、14、22 的平坦开阔接近 & 最高 2.1 m/s；同时要求距航点≥1.5
m、航向误差≤7°、横向偏差≤0.26 m、倾角≤13°、pitch≤6° \\
一般非快速路段 & 最高 1.2 m/s，再受地形状态约束 \\
WP26、28 & 0.95 m/s \\
WP16、24、25、27、29、31、32 & 0.70 m/s \\
WP26、27 的窄平台转弯 & 沿已通过来路后退 0.70 m，再转向 \\
WP28 最后一级台阶 & 配置 1.30 m 有界后退助跑 \\
WP28→29 & WP28 计分且倾角\textless12°持续0.4
s后，启用已知同层走廊；抑制跨层高度误判及 StepCommit，保留 LiDAR 与
BLOCKED/UNSTABLE/UNKNOWN \\
WP31、32 & 使用预设 body-clear 中间引导点；不推进官方航点 cursor \\
\end{longtable}

速度上限约束的是相应导航模式的命令，不等于实测速度或对所有恢复／引导模式统一生效的全局上限。StepCommit、route
hint 和控制交接另有命令限制。

\subsection{4.4
停滞检测与恢复}\label{ux505cux6edeux68c0ux6d4bux4e0eux6062ux590d}

Follower 在请求前进时并行监测：速度低于 0.08 m/s 持续 2.5
s；或到当前航点的最佳距离连续 12 s
没有有效改善。触发后进行有限倒退与转向，并继续尝试同一航点。StepCommit
有独立的推进超时和尝试预算。

需要区分 follower 的通用恢复与 Gate 16 router 的重试：\textbf{交付配置
{\ttfamily\small\seqsplit{max\_retries: 0}}}，Gate 16
失败不会自动获得多次额外攀爬机会。

\Needspace{240pt}

\section{5.
运动策略与部署接口约定}\label{ux8fd0ux52a8ux7b56ux7565ux4e0eux90e8ux7f72ux63a5ux53e3ux7ea6ux5b9a}

\subsection{5.1
模型分工与网络结构}\label{ux6a21ux578bux5206ux5de5ux4e0eux7f51ux7edcux7ed3ux6784}

\Needspace{6\baselineskip}
{\small\sffamily\bfseries\color{ReportGray}表 5. 运动策略网络与部署状态\par}
\vspace{-4pt}

\begin{longtable}[]{@{}
  >{\raggedright\arraybackslash}p{(\columnwidth - 4\tabcolsep) * \real{0.2000}}
  >{\raggedright\arraybackslash}p{(\columnwidth - 4\tabcolsep) * \real{0.5400}}
  >{\raggedright\arraybackslash}p{(\columnwidth - 4\tabcolsep) * \real{0.2600}}@{}}
\toprule\noalign{}
\begin{minipage}[b]{\linewidth}\raggedright
控制路径
\end{minipage} & \begin{minipage}[b]{\linewidth}\raggedright
结构
\end{minipage} & \begin{minipage}[b]{\linewidth}\raggedright
状态
\end{minipage} \\
\midrule\noalign{}
\endhead
\bottomrule\noalign{}
\endlastfoot
普通赛段 & 官方 57D 本体观测 actor & 使用官方权重和 runner \\
Gate 16 base & {\ttfamily\small\seqsplit{174→512→256→128→16}}，隐藏层
ELU & 冻结；来源 {\ttfamily\small\seqsplit{s10\_29cm\_stable.pt}} \\
Gate 16 residual &
{\ttfamily\small\seqsplit{174→512→512→256→128→16}}，隐藏层 ELU &
冻结部署；{\ttfamily\small\seqsplit{speed\_core}} iteration 120 \\
Residual critic & {\ttfamily\small\seqsplit{174→256→128→64→1}}，隐藏层
ELU & 仅训练使用 \\
连续楼梯 57D actor & 独立团队模型 & 打包但
{\ttfamily\small\seqsplit{stairs57\_enabled: false}} \\
\end{longtable}

两份 Gate 16 ONNX 的层尺寸与 ELU 算子已直接从模型文件核对，residual
checkpoint 的迭代、训练模式和 tensor shapes 也已读取。Actor 是
feed-forward MLP；上一动作作为显式输入，没有 RNN
或额外历史堆叠。{[}S4{]}

\subsection{5.2 174D 观测向量}\label{d-ux89c2ux6d4bux5411ux91cf}

\[
o_t=[0.25\omega_t^b,\ g_t^b,\ c_t,\ \tilde q_t-q_0,\ 0.05\dot q_t,\ a_{t-1},\ h_t]\in\mathbb R^{174}.
\]

\Needspace{6\baselineskip}
{\small\sffamily\bfseries\color{ReportGray}表 6. 174D 观测布局\par}
\vspace{-4pt}

\begin{longtable}[]{@{}
  >{\raggedright\arraybackslash}p{(\columnwidth - 4\tabcolsep) * \real{0.1600}}
  >{\raggedleft\arraybackslash}p{(\columnwidth - 4\tabcolsep) * \real{0.1000}}
  >{\raggedright\arraybackslash}p{(\columnwidth - 4\tabcolsep) * \real{0.7400}}@{}}
\toprule\noalign{}
\begin{minipage}[b]{\linewidth}\raggedright
切片
\end{minipage} & \begin{minipage}[b]{\linewidth}\raggedleft
维数
\end{minipage} & \begin{minipage}[b]{\linewidth}\raggedright
定义
\end{minipage} \\
\midrule\noalign{}
\endhead
\bottomrule\noalign{}
\endlastfoot
0:3 & 3 & 机身系角速度 ×0.25 \\
3:6 & 3 & 投影重力 \(R^{-1}[0,0,-1]\) \\
6:9 & 3 & 前向、横向和 yaw 速度命令 \\
9:25 & 16 & 关节位置相对默认值；轮角先置零 \\
25:41 & 16 & 包括轮速的关节速度 ×0.05 \\
41:57 & 16 & 上周期经过 guard 的最终 raw action \\
57:174 & 117 & 按训练约定转换的高度图 \\
\end{longtable}

\(\tilde q_t\) 为重排后将四个轮角置零的位置向量；\(c_t\) 为 runner
实际送入网络的速度条件，已包含适用的命令限幅。Policy order 为四组腿关节
FL、FR、HL、HR，随后四个轮子 FL、FR、HL、HR。默认角度为：

\par
\begin{minipage}{\linewidth}
\begin{ReportCode}
[0,-0.3,0.6, 0,-0.3,0.6, 0,0.3,-0.6, 0,0.3,-0.6, 0,0,0,0]
\end{ReportCode}
\end{minipage}
\par

该 observation 不包含实测机身线速度、接触标志、绝对航点或 router
phase；这些量可能供 router 或训练评测使用，但不是 actor 输入。

\subsection{5.3
基础策略与门控残差}\label{ux57faux7840ux7b56ux7565ux4e0eux95e8ux63a7ux6b8bux5dee}

令 \(b_t=\pi_b(o_t)\)、\(\delta_t=\mu_\theta(o_t)\)，部署计算为：

\[
a_t=\operatorname{clip}\left[b_t+m_t\,s\odot\operatorname{clip}(\delta_t,-4,4),-a_{\mathrm{max}},a_{\mathrm{max}}\right],
\]

其中
\(s=[1]^{12}\oplus[6]^4\)、\(a_{\mathrm{max}}=[8]^{12}\oplus[25]^4\)，\(m_t\in\{0,1\}\)
由 router arming、height skill gate 和 runner engagement
共同确定，\(a_{\mathrm{max}}\) 为逐通道 raw-action guard。

实际 skill gate 在中央列 2--6、行边界 4--8 检查向上高度差；连续两帧≥0.04
m 进入，至少保持 100 帧，连续 15 帧≤0.02 m 退出，最多保持 600
帧。保持阶段扩展到行边界
0--8。以上是\textbf{残差激活阈值}，不代表模型能通过任意该高度以上的障碍。

完整运行强制选择
{\ttfamily\small\seqsplit{gate16\_climb\_fallback}}：前向条件输入在低层限至
0.15 m/s，禁用 confidence fast adapter、policy mirroring 与 front-tuck
command profile。Bundle
中保留这些实验功能，不代表它们在该次运行中生效。{[}S10{]}

\subsection{5.4 动作解码与 PD
控制}\label{ux52a8ux4f5cux89e3ux7801ux4e0e-pd-ux63a7ux5236}

Raw action 重排到每条腿三个关节加一个轮子的 robot order：

\Needspace{6\baselineskip}
{\small\sffamily\bfseries\color{ReportGray}表 7. 动作解码参数\par}
\vspace{-4pt}

\begin{longtable}[]{@{}
  >{\raggedright\arraybackslash}p{(\columnwidth - 4\tabcolsep) * \real{0.1900}}
  >{\raggedright\arraybackslash}p{(\columnwidth - 4\tabcolsep) * \real{0.4300}}
  >{\raggedleft\arraybackslash}p{(\columnwidth - 4\tabcolsep) * \real{0.3800}}@{}}
\toprule\noalign{}
\begin{minipage}[b]{\linewidth}\raggedright
通道
\end{minipage} & \begin{minipage}[b]{\linewidth}\raggedright
目标
\end{minipage} & \begin{minipage}[b]{\linewidth}\raggedleft
缩放
\end{minipage} \\
\midrule\noalign{}
\endhead
\bottomrule\noalign{}
\endlastfoot
hip-x & 默认位置＋偏置 & 0.125 rad／raw unit \\
hip-y、knee & 默认位置＋偏置 & 0.25 rad／raw unit \\
wheel & 速度目标 & 5.0 rad/s／raw unit \\
\end{longtable}

腿使用 \(K_p=80,K_d=2\)，轮子使用 \(K_p=0,K_d=0.6\)，feed-forward torque
为零。SDK 输出为每执行器五列命令矩阵。

局部训练 runtime 每个策略步执行 20 个 1 ms 物理步，每步计算并裁剪 PD
torque：

\[
\tau=\operatorname{clip}\{K_p(q^*-q)+K_d(\dot q^*-\dot q),-\tau_{\mathrm{max}},\tau_{\mathrm{max}}\},
\]

其中腿关节 \(\tau_{\mathrm{max}}=50\) Nm，轮子为 14 Nm。Raw-action guard
与 torque clipping
分别约束不同变量，不能等同于完整的关节位置／速度及比赛扭矩合规证明。

\FloatBarrier

\section{6.
强化学习训练与模型选择}\label{ux5f3aux5316ux5b66ux4e60ux8badux7ec3ux4e0eux6a21ux578bux9009ux62e9}

\subsection{6.1
最终权重来源与训练阶段}\label{ux6700ux7ec8ux6743ux91cdux6765ux6e90ux4e0eux8badux7ec3ux9636ux6bb5}

实际使用的 residual checkpoint SHA-256 为
{\ttfamily\small\seqsplit{16e81160d099d345cb6766d06107a84f77ed02bb2ae9259e9f4076ece58eee74}}。Checkpoint
明确记录：

\par
\begin{minipage}{\linewidth}
\begin{ReportCode}
iteration = 120
terrain_mode = official_track
objective_mode = speed
activation_mode = heightmap
speed_range = [0.18, 0.40] m/s
yaw_range_rad = 0.1396263  (~8 degrees)
initial_residual_checkpoint = gate16_front_retention_20260817T121106/selected/best_checkpoint.pt
base_checkpoint = s10_29cm_stable.pt
\end{ReportCode}
\end{minipage}
\par

可确认的训练链为：既有 174D base 与前期 residual → Gate 16
成功轨迹修复和前轮保持研究 → {\ttfamily\small\seqsplit{speed\_core}}
速度目标 PPO → 固定矩阵选择 → 冻结并导出。早期材料描述了 57D→174D
warm-start、课程地形和 BC/DAgger，但本地没有完整保存 base
的训练日志及所有中间
selection，因此不能声称已经从原始初始化完整重建最终权重的训练历史。

最终训练 recipe 的关键文件在发布来源 {\ttfamily\small\seqsplit{5ef14fa}}
与所读本地训练仓库之间无差异；其中仍须区分 checkpoint 直接确认的设置与
launch script 提供的设置。{[}S4--S6{]}

\subsection{6.2
训练环境、重置分布与特权信号}\label{ux8badux7ec3ux73afux5883ux91cdux7f6eux5206ux5e03ux4e0eux7279ux6743ux4fe1ux53f7}

最终 {\ttfamily\small\seqsplit{speed\_core}} 使用
{\ttfamily\small\seqsplit{OfficialClosedLoopResidualEnv}}，直接加载官方
{\ttfamily\small\seqsplit{S10\_track.xml}}；固定台沿约 \(x=12.646\)
m，地面 \(z=0.10181425\) m，平台 \(z=0.47872480\) m，高差 0.37691055
m。环境最大 episode 为 850 个控制步，即 17 s。

训练入口采样如下：

\Needspace{6\baselineskip}
{\small\sffamily\bfseries\color{ReportGray}表 8. 训练入口分布\par}
\vspace{-4pt}

\begin{longtable}[]{@{}
  >{\raggedright\arraybackslash}p{(\columnwidth - 2\tabcolsep) * \real{0.3000}}
  >{\raggedright\arraybackslash}p{(\columnwidth - 2\tabcolsep) * \real{0.7000}}@{}}
\toprule\noalign{}
\begin{minipage}[b]{\linewidth}\raggedright
变量
\end{minipage} & \begin{minipage}[b]{\linewidth}\raggedright
{\ttfamily\small\seqsplit{speed\_core}} recipe
\end{minipage} \\
\midrule\noalign{}
\endhead
\bottomrule\noalign{}
\endlastfoot
入口距台沿 & 0.55--0.65 m；85\% 概率从
{\ttfamily\small\seqsplit{\{0.55,0.60,0.65\}}} 抽取，其余均匀采样 \\
命令前向速度 & 0.18--0.40 m/s \\
初始前向速度 & 与采样命令速度一致 \\
入口 yaw & 约 ±8° 均匀采样 \\
横向偏置 & ±0.02 m \\
机身高度扰动 & 此 recipe 为 0 \\
初始侧向速度、yaw rate & 默认 0 \\
地形 & 固定官方 Gate 16 \\
\end{longtable}

Checkpoint 中保存的
{\ttfamily\small\seqsplit{depth\_range={[}0.3,0.4{]}}} 是 trainer
通用字段；{\ttfamily\small\seqsplit{terrain\_mode=official\_track}}
不使用可变高度 fixture，\textbf{不能据此声称该阶段进行了 30--40 cm
台阶随机化}。本阶段也没有证据支持质量、摩擦、驱动延迟或 LiDAR
噪声随机化。

局部环境在控制周期内还设置 yaw 命令为
\(\operatorname{clip}(-1.5\psi,-0.8,0.8)\)，除非显式指定。因而模型训练并非在完全无辅助的任务导航下进行。

部署 fallback 的网络前向速度条件最高为 0.15
m/s，低于本阶段训练采样的下界 0.18 m/s；入口距离 0.62--0.70 m
也部分超出训练的 0.55--0.65 m。较低速度出现在前期 recipe
和发布选择矩阵中，但当前证据不能证明训练与部署入口分布完全覆盖。

Actor 与 critic 均接收 174D observation；critic 没有额外 privileged
observation。但奖励和 episode
结束条件使用仿真轮心、平台高度与机身状态，属于 privileged
reward/evaluation signals。

\subsection{6.3 PPO
优化目标与实现}\label{ppo-ux4f18ux5316ux76eeux6807ux4e0eux5b9eux73b0}

Residual actor 在训练中定义独立高斯分布：

\[
\delta_t\sim\mathcal N(\mu_\theta(o_t),\operatorname{diag}(\sigma^2)),
\]

策略评估和 ONNX 部署使用均值。Base actor 冻结且不进入
optimizer；residual actor 从已选 checkpoint 初始化，critic
在切换到速度目标时重新初始化。

采用 GAE，\(\gamma=0.99,\lambda=0.95\)，batch 内标准化 advantage。PPO
ratio 基于 clipping/gating 前采样的 residual
action：\(\rho_t=\pi_\theta(\delta_t\mid o_t)/\pi_{\theta_{\mathrm{old}}}(\delta_t\mid o_t)\)，\(\hat R_t=\hat A_t+V_{\mathrm{old}}(o_t)\)。Actor
解冻后的联合损失为：

\[
\begin{aligned}
\mathcal L={}&-\mathbb E\left[\min\left(\rho_t\hat A_t,\operatorname{clip}(\rho_t,1-\epsilon,1+\epsilon)\hat A_t\right)\right]\\
&+0.5\mathbb E\left[(V(o_t)-\hat R_t)^2\right]-\beta\mathcal H(\pi_\theta)\\
&+\alpha\mathbb E\left[\frac{1}{16}\sum_{j=1}^{16}(\mu_{\theta,j}(o_t)-\mu_{\mathrm{ref},j}(o_t))^2\right].
\end{aligned}
\]

\(\mu_{\mathrm{ref}}\) 为该阶段初始化后冻结的 residual actor；参数名虽为
{\ttfamily\small\seqsplit{zero\_anchor\_coef}}，在 warm-start
情况下约束目标是初始策略输出，不是零动作。

实现将成功、跌倒、动作发散及 850 步时间上限统一返回为
{\ttfamily\small\seqsplit{done}}，GAE 对这些 transition 均令
{\ttfamily\small\seqsplit{nonterminal=0}}，不为时间上限额外
bootstrap；仅 rollout batch 截断且 episode
尚未结束时保留下一状态价值。前 40 次 critic warm-up 仅更新 value
loss，actor、entropy 与 reference-anchor 项不参与更新。

\Needspace{6\baselineskip}
{\small\sffamily\bfseries\color{ReportGray}表 9. 最终 PPO 阶段超参数\par}
\vspace{-4pt}

\begin{longtable}[]{@{}
  >{\raggedright\arraybackslash}p{(\columnwidth - 4\tabcolsep) * \real{0.3700}}
  >{\raggedright\arraybackslash}p{(\columnwidth - 4\tabcolsep) * \real{0.3500}}
  >{\raggedright\arraybackslash}p{(\columnwidth - 4\tabcolsep) * \real{0.2800}}@{}}
\toprule\noalign{}
\begin{minipage}[b]{\linewidth}\raggedright
参数
\end{minipage} & \begin{minipage}[b]{\linewidth}\raggedright
{\ttfamily\small\seqsplit{speed\_core}} 设置
\end{minipage} & \begin{minipage}[b]{\linewidth}\raggedright
依据
\end{minipage} \\
\midrule\noalign{}
\endhead
\bottomrule\noalign{}
\endlastfoot
Residual actor hidden & 512, 512, 256, 128，ELU & Checkpoint＋ONNX \\
Critic hidden & 256, 128, 64，ELU & Checkpoint＋trainer \\
环境数／每环境 rollout & 12／32 & Launch script \\
每迭代 transition 数 & 384 & 上述配置计算 \\
Optimizer & Adam & Trainer \\
Actor／critic LR & \(2\times10^{-8}\)／\(1\times10^{-5}\) & Launch
script \\
Epochs／minibatches & 1／4 & Launch script \\
PPO clip \(\epsilon\) & 0.05 & Launch script \\
Entropy coefficient \(\beta\) & 0 & Launch script \\
初始 log standard deviation & −3.5 & Launch script \\
Reference-anchor coefficient \(\alpha\) & 2.0 & Launch script \\
Critic warm-up & 前 40 次迭代冻结 actor & Launch script＋trainer \\
Gradient-norm clipping & 1.0 & PPO implementation \\
log-std bounds & {[}−5,0{]} & PPO implementation \\
评估／保存间隔 & 40 次迭代 & Launch script \\
计划最大迭代数 & 200；连续3次评估不改善可停止 & Launch script \\
实际部署 checkpoint & iteration 120 & Checkpoint \\
训练 seed & 377511 & Launch script \\
\end{longtable}

按该 recipe，iteration 120 对应本阶段 46,080 个 rollout
transitions，不含先前训练及评估交互；这是配置推算值，不能当作总训练样本量。12
个环境由单个 env-step worker 串行推进，神经网络使用批量推理／CUDA
更新，未使用 Isaac Gym 式 GPU 并行物理仿真。当前资料不足以报告端到端
GPU-hours 或整个训练链的样本效率。

\subsection{6.4
速度导向奖励函数}\label{ux901fux5ea6ux5bfcux5411ux5956ux52b1ux51fdux6570}

令 \(\Delta x\) 为裁剪到 {[}−0.05,0.05{]} m
的单步前进量；\(\Delta F,\Delta R,\Delta C\)
为前轮、后轮及机身进展的历史最优值增量；\(\Delta N\)
为历史最多越沿轮数的增量。Residual 激活时的稠密奖励为：

\[
\begin{aligned}
r_t^{\mathrm{active}}={}&35\Delta x+20\Delta F+70\Delta R+35\Delta C+6\Delta N-0.05\\
&-0.0005\operatorname{mean}(\bar\delta_t^2)\\
&-0.0005\operatorname{mean}[(\bar\delta_t-\bar\delta_{t-1})^2].
\end{aligned}
\]

\(\bar\delta\) 是 clipping/gating 后、乘各通道 scale 前的 residual。令
\([u]_0^1=\operatorname{clip}(u,0,1)\)、\(H=z_{\mathrm{deck}}-z_{\mathrm{floor}}\)，前后轮瞬时进展分别为：

\[
\begin{aligned}
F_t^{\mathrm{inst}}&=\frac12\sum_{i\in\mathrm{front}}\left[\frac{x_i-x_{\mathrm{lip}}+0.45}{0.35}\right]_0^1\left[\frac{z_i-z_{\mathrm{floor}}-r_w}{H}\right]_0^1,\\
R_t^{\mathrm{inst}}&=\frac12\sum_{i\in\mathrm{rear}}\left[\frac{x_i-x_{\mathrm{lip}}+0.35}{0.35}\right]_0^1\left[\frac{z_i-z_{\mathrm{floor}}-r_w}{H}\right]_0^1.
\end{aligned}
\]

机身进展为
\(C_t^{\mathrm{inst}}=[(x_{\mathrm{base}}-x_{\mathrm{lip}}+0.35)/0.70]_0^1\)。各项先更新历史最大值，再取相对前一步历史最大值的增量。代码字段
{\ttfamily\small\seqsplit{best\_com\_progress}} 使用 base
position，不是质量加权全身质心。

完整奖励为激活指示量乘上述稠密奖励，再独立累加成功、跌倒与发散项：

\[
\begin{aligned}
r_t={}&\mathbf1_{\mathrm{active}}r_t^{\mathrm{active}}
+\mathbf1_{\mathrm{success}}[300+0.5\max(0,850-n_t)]\\
&-250\mathbf1_{\mathrm{fallen}}-300\mathbf1_{\mathrm{diverged}},
\end{aligned}
\]

其中 \(n_t\) 为当前 episode 已执行的控制步数。事件奖励不要求 residual
激活；成功与跌倒标志独立计算，代码没有互斥优先级。所有历史最佳增量非负，前轮重新掉下不会自动撤销既得进度。{\ttfamily\small\seqsplit{objective\_mode=speed}}
不加入前轮保持或收腿阶段奖励。

局部环境的成功条件是四个轮心均满足 \(x_i\ge x_{\mathrm{lip}}+0.02\) 且
\(z_i\ge z_{\mathrm{deck}}+0.70r_w\)，\(r_w=0.081\)
m。跌倒条件为投影重力 z 分量\textgreater−0.20，或机身高度低于 floor+0.12
m。该成功条件没有部署 router
的接触数量、24°姿态约束和持续确认，因此局部成功率不能直接代替整圈
handoff 成功率。

\subsection{6.5
课程训练、行为克隆与失败状态聚合}\label{ux8bfeux7a0bux8badux7ec3ux884cux4e3aux514bux9686ux4e0eux5931ux8d25ux72b6ux6001ux805aux5408}

训练仓库实现了以下开发路径，实际采用情况由发布记录区分：

\begin{itemize}
\tightlist
\item
  \textbf{57D→174D warm-start：} 保留第一层原 57 列权重，将新增 117
  列置零；早期 recipe 包含 flat、8 cm、14 cm、20 cm
  地形阶段。它是模型扩展方法和早期训练配置，不足以证明当前 base
  完整经历这些阶段。
\item
  \textbf{成功轨迹监督：} 收集 CEM
  和已有策略的四轮成功轨迹，将最终动作与固定 base 输出之差，按 residual
  scale 还原为监督目标；关键接触阶段可重采样。
\item
  \textbf{Failure-state aggregation：} 运行当前 residual，收集失败
  rollout 的相关状态，以成功轨迹中的近邻提供 pseudo-target，并与成功样本
  rehearsal、参数锚定共同优化。后期版本将修正限制在
  {\ttfamily\small\seqsplit{settle/push}}
  状态。该方法的标签来自近邻成功状态，不能等同于每个失败状态都有可靠的在线专家动作。
\item
  \textbf{Front-retention curriculum：} 脚本设置
  0.10--0.12、0.10--0.20、0.10--0.30 m/s 三阶段，分别计划 180、260、360
  次 PPO 迭代，逐步扩大 yaw。最终
  {\ttfamily\small\seqsplit{speed\_core}} checkpoint 指向这一阶段输出的
  selected
  checkpoint，但缺少完整中间材料来确认每阶段实际执行次数与获选模型。
\item
  \textbf{Speed fine-tuning：} 最终采用上述
  {\ttfamily\small\seqsplit{speed\_core}}。后续 phase-reward PPO、DAgger
  小步更新及 checkpoint blending
  未通过既定可靠性筛选，发布模型保持冻结。
\end{itemize}

因此，不能将``所有研究分支依次成功训练并共同提升最终模型''作为本项目的方法描述。{[}S5--S7{]}

\subsection{6.6
模型选择与确定性评测}\label{ux6a21ux578bux9009ux62e9ux4e0eux786eux5b9aux6027ux8bc4ux6d4b}

正式局部选择矩阵包含 45 个状态：距离
{\ttfamily\small\seqsplit{\{0.55,0.60,0.65\}}} m × yaw
{\ttfamily\small\seqsplit{\{−12°,0°,12°\}}} × 速度
{\ttfamily\small\seqsplit{\{0.10,0.15,0.20,0.25,0.30\}}} m/s。另设 27
个速度筛选状态，距离相同，yaw
{\ttfamily\small\seqsplit{\{−8°,0°,8°\}}}，速度
{\ttfamily\small\seqsplit{\{0.20,0.30,0.40\}}} m/s。

Speed selection 先要求两个矩阵的成功率均不低于相应 baseline，再比较
failure-penalized completion steps：成功使用实际步数，所有失败统一计为
850 步。后期 release selection 另约束 fall rate 和 drop-free success
不退化。

训练中的周期评估恢复各环境 RNG state，保持 common random
numbers；发布评估保留固定 case 顺序，相关 DAgger 校验使用 CPU actor
和每个速度复用一个环境，以减少数值与环境生命周期差异。固定矩阵参与了选模，属于
validation/selection set，不是独立无偏 test set。

\Needspace{165pt}

\subsection{6.7 ONNX
导出与模型身份}\label{onnx-ux5bfcux51faux4e0eux6a21ux578bux8eabux4efd}

导出仅包含 actor mean，base 与 residual 分别保存。发布清单记录的
PyTorch→ONNX 最大绝对误差分别为 \(3.814697\times10^{-6}\) 和
\(4.768372\times10^{-6}\)；本次文档核查验证了文件 SHA-256
与结构，未重新执行数值推理比对。

\Needspace{6\baselineskip}
{\small\sffamily\bfseries\color{ReportGray}表 10. 部署模型 SHA-256\par}
\vspace{-4pt}

\begin{longtable}[]{@{}
  >{\raggedright\arraybackslash}p{(\columnwidth - 2\tabcolsep) * \real{0.2700}}
  >{\raggedright\arraybackslash}p{(\columnwidth - 2\tabcolsep) * \real{0.7300}}@{}}
\toprule\noalign{}
\begin{minipage}[b]{\linewidth}\raggedright
文件
\end{minipage} & \begin{minipage}[b]{\linewidth}\raggedright
SHA-256
\end{minipage} \\
\midrule\noalign{}
\endhead
\bottomrule\noalign{}
\endlastfoot
{\ttfamily\small\seqsplit{policy.onnx}} &
{\ttfamily\small\seqsplit{5c1b388f951b282693b4497cd4fd2fd1b53fe1bcd989758c0af42f140a20fb16}} \\
{\ttfamily\small\seqsplit{climb\_residual.onnx}} &
{\ttfamily\small\seqsplit{de61441facb21f0301787b26f168f8447fdc0d83c337eeea884537bbf2468889}} \\
\end{longtable}

模型接入来源为
{\ttfamily\small\seqsplit{belsun/goai-s10-gate16-policy@216b77a}}，asset
bundle 为 {\ttfamily\small\seqsplit{b6535a4}}。模型权重身份与
runner/config 身份都影响结果，不能只比较 ONNX hash
就认定两个实验配置相同。

\Needspace{500pt}

\section{7.
规则策略路由与控制器交接}\label{ux89c4ux5219ux7b56ux7565ux8defux7531ux4e0eux63a7ux5236ux5668ux4ea4ux63a5}

\subsection{7.1
状态机与入口条件}\label{ux72b6ux6001ux673aux4e0eux5165ux53e3ux6761ux4ef6}

\begin{center}
\includegraphics[width=\linewidth,height=0.52\textheight,keepaspectratio]{report_assets/figure-2.pdf}
\captionof{figure}{Gate 16 状态机与主要转移条件；交付配置的重试预算为 0。}
\end{center}

Gate 16 的边缘中心、法向和上平台高度来自预先测量的 MJCF
配置；当前不是通用在线台沿重建。稳定接管要求：

\Needspace{185pt}
{\small\sffamily\bfseries\color{ReportGray}表 11. Gate 16 接管入口条件\par}
\vspace{-4pt}

\begin{longtable}[]{@{}
  >{\raggedright\arraybackslash}p{(\columnwidth - 2\tabcolsep) * \real{0.5600}}
  >{\raggedleft\arraybackslash}p{(\columnwidth - 2\tabcolsep) * \real{0.4400}}@{}}
\toprule\noalign{}
\begin{minipage}[b]{\linewidth}\raggedright
条件
\end{minipage} & \begin{minipage}[b]{\linewidth}\raggedleft
当前 fallback
\end{minipage} \\
\midrule\noalign{}
\endhead
\bottomrule\noalign{}
\endlastfoot
沿台沿法向的机身距离 & 0.62--0.70 m \\
实测前向速度 & 0.08--0.20 m/s \\
目标速度命令 & 0.18 m/s；runner 中最高 0.15 m/s \\
横向位置误差 & ≤0.25 m \\
朝向误差 & ≤6° \\
yaw rate & ≤0.10 rad/s \\
组合倾角 & ≤12° \\
条件持续时间 & 0.10 s \\
\end{longtable}

官方 actor 拥有移动接近阶段的控制权。Gate 16 可提前 shadow inference，但
shadow
输出不能降低接近速度或写入实际动作历史。接管边沿从实测关节位置和轮速逆解上一动作，避免使用尚未执行的
shadow action 或另一策略不兼容的历史。

\subsection{7.2
四轮净空与受控退出}\label{ux56dbux8f6eux51c0ux7a7aux4e0eux53d7ux63a7ux9000ux51fa}

部署 verification 要求：四轮心沿冻结台沿法向均超过边缘 0.02
m；四轮心高度≥0.4787248+0.70×0.081=0.5354248 m；至少三个轮子有
terrain-contact flag；倾角≤24°；进入确认时平面速度≤1.08
m/s。组合条件持续 0.10 s，允许 0.30 s unready chatter，确认超时为 5.0
s。

通过后依次撤销 residual arming、释放 Gate 16 owner、进入 0.25 s
measured-position safe hold、重置官方 actor 内部历史、等待
{\ttfamily\small\seqsplit{/joints/owner=official}}、清空 follower
transient state，并以 0.50 m/s 恢复 0.80 s。

移动接管是 arbiter 的受限例外：同一 SDK tick 已有有效 Gate 16 command
matrix 才可直接切入；否则短暂保留官方输出，超过 0.10 s 未准备好则
hold。控制权互斥由 SDK 输出边界保证，router
的请求字符串本身不等于执行器已交接。

\FloatBarrier

\section{8.
时序、保护机制与失败语义}\label{ux65f6ux5e8fux4fddux62a4ux673aux5236ux4e0eux5931ux8d25ux8bedux4e49}

\subsection{8.1
时间与更新频率}\label{ux65f6ux95f4ux4e0eux66f4ux65b0ux9891ux7387}

局部 RL 环境严格按 20 次×1 ms 物理积分执行一个 20 ms policy step。完整
ROS/SDK 系统各模块名义上以 50 Hz
更新，但进程调度、消息到达和仿真推进不是同一件事。

在本文固定源码中，router 使用 ROS node clock，部分 dwell 通过固定
{\ttfamily\small\seqsplit{dt=1/control\_rate}} 累加；Gate 16 C++ runner
使用 monotonic clock。官方路线计时来自 MuJoCo
{\ttfamily\small\seqsplit{sim\_time}}。因此，\textbf{本报告不宣称所有
watchdog、dwell 和控制切换均已统一到
{\ttfamily\small\seqsplit{/clock}}，也不将名义 50 Hz
当作实测实时性保证}。计算负载造成的 callback/cadence
变化可影响控制结果。

Router 的 odometry、LiDAR 和 heightmap freshness 时间记录为 callback
接收时刻，而非统一的传感器采样时刻；轮子状态也没有独立 age 字段。现有
watchdog
主要检测消息停止到达，不能完整识别排队中的陈旧样本或轮状态单独停更。

\Needspace{260pt}

\subsection{8.2
实际保护条件}\label{ux5b9eux9645ux4fddux62a4ux6761ux4ef6}

\Needspace{6\baselineskip}
{\small\sffamily\bfseries\color{ReportGray}表 12. 运行保护与退出条件\par}
\vspace{-4pt}

\begin{longtable}[]{@{}
  >{\raggedright\arraybackslash}p{(\columnwidth - 2\tabcolsep) * \real{0.4500}}
  >{\raggedright\arraybackslash}p{(\columnwidth - 2\tabcolsep) * \real{0.5500}}@{}}
\toprule\noalign{}
\begin{minipage}[b]{\linewidth}\raggedright
事件
\end{minipage} & \begin{minipage}[b]{\linewidth}\raggedright
处理
\end{minipage} \\
\midrule\noalign{}
\endhead
\bottomrule\noalign{}
\endlastfoot
初始所需观测未齐 & 零速度，保留官方 ownership \\
运行时观测 age\textgreater0.5 s & 停止；攀爬中取消策略并转恢复 \\
持续 stale 状态达到3.0 s & Abort \\
倾角≥60° & Abort \\
Gate 16 无有效进展 & 12 s progress window 后恢复／退出 \\
Gate 16 攀爬超时 & 20 s \\
Alignment 超时 & 40 s \\
非有限输出、尺寸错误、失效命令 & 拒绝并 safe hold \\
请求 owner 与实际 owner 不一致 & 等待确认，不假定交接已完成 \\
\end{longtable}

代码在 \textbf{Gate 16/stairs 委托控制处于
{\ttfamily\small\seqsplit{CLIMB/VERIFY\_CLEAR}} 时}检查实测关节
torque；任一关节绝对值\textgreater50 Nm 连续1.0
s触发退出。普通导航阶段清空该 watchdog
的累计状态。它使用全关节统一阈值，不能替代轮子14
Nm的逐执行器限幅或比赛独立的持续过扭矩规则。局部训练的每步 torque
clipping、上游执行器边界和运行期 watchdog
是不同层次，本次完赛记录不构成完整 torque-duration 合规证明。

\FloatBarrier

\section{9.
实验协议与结果}\label{ux5b9eux9a8cux534fux8baeux4e0eux7ed3ux679c}

\subsection{9.1
连续系统自测}\label{ux8fdeux7eedux7cfbux7edfux81eaux6d4b}

使用 {\ttfamily\small\seqsplit{3660b81}} 代码和固定上游，独立 ARM64
Docker 构建；segment harness 仅提供 WP0
初始状态与观测记录，之后连续运行生产控制栈。该结果是 team
self-test，未被表述为主办方认证成绩。

\Needspace{6\baselineskip}
{\small\sffamily\bfseries\color{ReportGray}表 13. 完整路线自测结果\par}
\vspace{-4pt}

\begin{longtable}[]{@{}
  >{\raggedright\arraybackslash}p{(\columnwidth - 2\tabcolsep) * \real{0.4400}}
  >{\raggedleft\arraybackslash}p{(\columnwidth - 2\tabcolsep) * \real{0.5600}}@{}}
\toprule\noalign{}
\begin{minipage}[b]{\linewidth}\raggedright
指标
\end{minipage} & \begin{minipage}[b]{\linewidth}\raggedleft
结果
\end{minipage} \\
\midrule\noalign{}
\endhead
\bottomrule\noalign{}
\endlastfoot
Seed & 8 \\
按序航点事件 & 33／33，WP0--WP32 \\
WP0／WP32 仿真时间 & 0.001／392.258 s \\
路线 elapsed & \textbf{392.257 s} \\
Recorder wall elapsed & 674.01 s \\
行驶距离 & 249.94 m \\
最终 WP32 距离 & 0.177 m \\
最大倾角 & 44.6° \\
Recorder stalls & 2 \\
漏点／跌倒／超时／router abort & 均未发生 \\
\end{longtable}

WP16 在 152.871 s 通过，WP17 在 155.488 s 通过。日志保留 fallback
接管、residual arm、四轮验证、safe hold、官方 owner 确认及 follower
reset。WP28→29 用时 8.367 s，该段启用了同层走廊修正。

Recorder 在严格 0.18 m 终点条件满足后结束 launch，未必等待后续 router
{\ttfamily\small\seqsplit{DONE}} 状态锁存；实际最后 owner 为
official。完成判据来自独立官方航点序列和 timer，不能以进程 exit code
或单个 {\ttfamily\small\seqsplit{reached=true}} 字段替代。{[}S1--S2{]}

\Needspace{260pt}

\subsection{9.2
局部策略结果与适用范围}\label{ux5c40ux90e8ux7b56ux7565ux7ed3ux679cux4e0eux9002ux7528ux8303ux56f4}

\Needspace{6\baselineskip}
{\small\sffamily\bfseries\color{ReportGray}表 14. 局部策略评测结果\par}
\vspace{-4pt}

\begin{longtable}[]{@{}
  >{\raggedright\arraybackslash}p{(\columnwidth - 6\tabcolsep) * \real{0.1800}}
  >{\raggedright\arraybackslash}p{(\columnwidth - 6\tabcolsep) * \real{0.2500}}
  >{\raggedright\arraybackslash}p{(\columnwidth - 6\tabcolsep) * \real{0.2900}}
  >{\raggedright\arraybackslash}p{(\columnwidth - 6\tabcolsep) * \real{0.2800}}@{}}
\toprule\noalign{}
\begin{minipage}[b]{\linewidth}\raggedright
实验
\end{minipage} & \begin{minipage}[b]{\linewidth}\raggedright
配置／样本
\end{minipage} & \begin{minipage}[b]{\linewidth}\raggedright
结果
\end{minipage} & \begin{minipage}[b]{\linewidth}\raggedright
解释
\end{minipage} \\
\midrule\noalign{}
\endhead
\bottomrule\noalign{}
\endlastfoot
{\ttfamily\small\seqsplit{speed\_core}} 发布矩阵 & 45 个固定入口状态 &
38/45 成功；2/45 跌倒；20/45 drop-free 成功 &
固定台阶上的候选选择证据 \\
同一矩阵成功样本用时 & 38 个成功 case & 平均 6.134 s；总体标准差 3.098 s
& 从 episode reset 起计，不含失败样本 \\
同一矩阵失败惩罚后均值 & 45 个 case，失败计17 s & 7.824 s &
与成功样本均值不同 \\
同一矩阵前轮→后轮完成 & 38 个成功 case & 平均208.26步，约4.165 s &
从双前轮首次满足几何条件到四轮同时满足 \\
后续 front-tuck/profile 实验 & 独立局部适配配置 &
发布文档报告39/45；固定距离速度的斜入扫描25/25 &
本次完整运行未启用，不计入部署鲁棒性 \\
\end{longtable}

{\ttfamily\small\seqsplit{drop-free success}}
指最终成功，且双前轮首次满足 §6.4
的轮心几何条件后、四轮完成前，未出现满足条件的前轮数从2降到小于2的事件。该指标的``支撑''是几何代理，未检查接触力；20/45
的分母为全部 case。表中标准差按总体标准差计算，未作 \(N-1\)
校正。{[}S11{]}

Checkpoint 内嵌 {\ttfamily\small\seqsplit{eval\_success=0.6667}}
来自该训练阶段的内部评估分布，与正式45-case矩阵不同。它和38/45不能混用，二者也都不等于全路线成功率。{[}S4、S7{]}

\subsection{9.3
负面证据与当前评测缺口}\label{ux8d1fux9762ux8bc1ux636eux4e0eux5f53ux524dux8bc4ux6d4bux7f3aux53e3}

保留测试中存在 Gate 16
跨沿后无法交接、其他楼梯段倾倒和感知更新中断。另一次独立 Mac
批次五次均失败，其中三次为启动感知问题、两次为途中倾倒；其运行条件不应与392秒成功样本直接合并成控制算法成功率。{[}S8{]}

目前证据不足以给出统一协议下的整圈成功率、置信区间或相对 baseline
的因果增益。尚缺：独立 held-out 入口／地形测试、固定硬件负载下多 seed
整圈统计、逐模块消融、策略切换
baseline，以及真机重复实验。现有课程特例和选模矩阵使用也限制了泛化结论。392.257
s
是保留并用于提交的成功样本，不是全部试验的平均成绩，也未据此主张当前所有版本的最佳成绩。

\subsection{9.4
代码验证与回放}\label{ux4ee3ux7801ux9a8cux8bc1ux4e0eux56deux653e}

提交包报告：非 joint-owner fixture 的源码测试 378/378
通过，training/observation contract 测试 29/29 通过；另有未启用 adaptive
owner 路径的 C++ fixture 检查失败，产生9个 pytest setup
errors。因此整体测试状态不是全绿。Fallback owner
生命周期在该次完整运行中得到验证，但不能用单次运行消除 fixture 缺口。

720p 回放由同一次状态记录离线渲染，30 fps、11,768帧、392.267 s，与官方
elapsed 差0.010
s。该视频是仿真状态重绘。以上实验与测试均为归档结果；本次文档修订未重新执行仿真或测试套件。

\Needspace{240pt}

\section{10.
部署与可复现性}\label{ux90e8ux7f72ux4e0eux53efux590dux73b0ux6027}

\subsection{10.1 运行环境}\label{ux8fd0ux884cux73afux5883}

Ubuntu 24.04、ROS 2 Jazzy、Python 3.12；Docker base image 按 digest
固定，Python dependencies 有 lock file。提交包记录 NumPy 1.26.4、MuJoCo
3.11.0、ONNX Runtime 1.28.0。比赛运行无商业
API、托管模型或在线推理依赖。

在正确 checkout 和已授权的上游资源就绪后，主要流程为：

\par
\begin{minipage}{\linewidth}
\begin{ReportCode}
S10_UPSTREAM_OFFLINE=1 docker compose run --rm s10 scripts/setup_upstream.sh
docker compose run --rm s10 scripts/build.sh
docker compose run --rm s10 scripts/verify_install.sh
docker compose run --rm s10 scripts/run_race.sh --headless
\end{ReportCode}
\end{minipage}
\par

{\ttfamily\small\seqsplit{run\_race.sh}} 启用比赛 router 配置；底层
launch 的 opt-in 默认值及旧注释不能代替实际入口脚本。独立实验应隔离 ROS
domain、容器和 build volume，并在运行前校验安装源码与模型身份。

完整 seed-8 复测入口为：

\par
\begin{minipage}{\linewidth}
\begin{ReportCode}
docker compose run --rm s10 scripts/run_segment.sh \
  --start 0 --end 32 --seeds 1 --seed-from 8 --max-time 2400 \
  --out /ws/results/retest_seed8 --tag full --router \
  --router-params /ws/src/s10_bringup/config/strategy_gate16.yaml
\end{ReportCode}
\end{minipage}
\par

这提供复测流程，不保证另一计算负载或机器上重现完全相同的接触轨迹和完成时间。

\subsection{10.2
训练复现条件}\label{ux8badux7ec3ux590dux73b0ux6761ux4ef6}

{\ttfamily\small\seqsplit{training/run\_gate16\_speed\_first\_pipeline.sh}}
接收上一阶段的 selected run 与输出目录，执行 speed matrix baseline、PPO
候选、矩阵筛选和双 ONNX 导出。运行需要官方
XML、{\ttfamily\small\seqsplit{s10\_29cm\_stable.pt}}、前期
residual、前期 {\ttfamily\small\seqsplit{selection.json}}
及其引用的评测摘要。

当前本地保留最终 residual checkpoint、两份
ONNX、发布选择摘要和训练脚本；缺少 base 原始训练日志、完整前期
checkpoint／dataset 链及 {\ttfamily\small\seqsplit{speed\_core}}
原始训练曲线。\textbf{部署复测材料较完整，最终模型从头再训练的材料仍不完整。}
本文据此报告精确可核对的算法与
recipe，不给出未经证据支持的收敛曲线、总训练成本或所有课程阶段成绩。

\subsection{10.3
模型与依赖归属}\label{ux6a21ux578bux4e0eux4f9dux8d56ux5f52ux5c5e}

普通 actor、机器人资产和官方环境来自主办方；Gate 16 base/residual
和实验楼梯模型属于团队贡献。当前归档披露部分模型 bundle
缺少独立许可证，公开分发前仍需补齐贡献者授权与归属；官方受限资源应按其授权渠道提供。本文不将第三方能力计为团队从零实现的算法。

\FloatBarrier

\section{11.
技术局限与后续研究}\label{ux6280ux672fux5c40ux9650ux4e0eux540eux7eedux7814ux7a76}

\begin{enumerate}
\def\labelenumi{\arabic{enumi}.}
\tightlist
\item
  \textbf{Perception/localization transfer：}
  用物理传感器重建高度图，替换真值定位和轮子接触 oracle，加入
  validity、协方差与时间同步。
\item
  \textbf{Training--deployment mismatch：} 最终训练与 fallback
  部署的速度条件、入口距离、完成判据和控制节拍不完全一致；需要在完整部署入口包络内验证，并对齐时间与终止语义。
\item
  \textbf{Robustness：} 最终速度阶段只随机化有限入口状态，没有完整
  dynamics randomization；规则的放行窗口也没有得到连续状态空间验证。
\item
  \textbf{Course dependence：} 固定台沿、速度白名单、平台特例及 route
  hints 对已知赛道有效，尚无跨场景成功率。
\item
  \textbf{Evaluation completeness：} 补充重复整圈、held-out
  测试、机制消融、真实推理延迟／控制 jitter 和力矩时间序列分析。
\item
  \textbf{Decision learning：} Learned router
  可作为后续研究方向，但需以当前规则系统为
  baseline，并保留单一控制权及异常保护；该方向尚无实现或实验结果。
\end{enumerate}

\clearpage

\section{12.
源码与证据索引}\label{ux6e90ux7801ux4e0eux8bc1ux636eux7d22ux5f15}

正文中的实现与数字来自以下材料。局部训练代码与部署代码分别索引，避免将训练
runner、实验 adapter 和比赛配置混为一体。本次已将 router、ROS
adapter、pursuit、导航／策略配置、policy runner、joint
arbiter、heightmap sampler 和启动脚本等九个关键部署文件与 Git
{\ttfamily\small\seqsplit{3660b81}} 逐字节比较，均一致。

\begin{itemize}
\tightlist
\item
  \textbf{S1 --- 提交技术设计与验证报告：}
  \href{<workspace-legacy>/resources/submission_update_20260820_392s/submission_20260820_392s_main3660b81/03_documents/TECHNICAL_DESIGN.md}{TECHNICAL\_DESIGN.md}、\href{<workspace-legacy>/resources/submission_update_20260820_392s/submission_20260820_392s_main3660b81/04_evidence/validation_report.md}{validation\_report.md}。遇到文档与代码不一致，本报告采用实际代码与运行记录。
\item
  \textbf{S2 --- 连续自测证据：}
  \href{<workspace-legacy>/resources/submission_update_20260820_392s/submission_20260820_392s_main3660b81/04_evidence/RUN_SUMMARY.md}{RUN\_SUMMARY.md}、\href{<workspace-legacy>/resources/submission_update_20260820_392s/submission_20260820_392s_main3660b81/04_evidence/raw/00_32_seed8.log}{原始日志}。
\item
  \textbf{S3 --- 部署源码：}
  \href{<workspace-legacy>/resources/submission_mac_retest_20260820/20260820_215224_main3660b81/source/goai26-s10-racing/src/s10_auto_nav/s10_auto_nav/follower_node.py}{follower\_node.py}、\href{<workspace-legacy>/resources/submission_mac_retest_20260820/20260820_215224_main3660b81/source/goai26-s10-racing/src/s10_auto_nav/s10_auto_nav/strategy/router.py}{router.py}、\href{<workspace-legacy>/resources/submission_mac_retest_20260820/20260820_215224_main3660b81/source/goai26-s10-racing/src/s10_bringup/config/nav.yaml}{nav.yaml}、\href{<workspace-legacy>/resources/submission_mac_retest_20260820/20260820_215224_main3660b81/source/goai26-s10-racing/src/s10_bringup/config/strategy_gate16.yaml}{strategy\_gate16.yaml}。
\item
  \textbf{S4 --- 本次直接模型核查：}
  \href{<workspace>/TECHNICAL_MODEL_AUDIT.json}{TECHNICAL\_MODEL\_AUDIT.json}，包含
  checkpoint metadata、tensor shapes、ONNX 算子和 SHA-256。Metadata
  读取不执行训练代码。
\item
  \textbf{S5 --- 最终阶段 recipe 与 PPO：}
  \href{<workspace-legacy>/goai-s10-gate16-policy-v1-5/training/run_gate16_speed_first_pipeline.sh}{run\_gate16\_speed\_first\_pipeline.sh}、\href{<workspace-legacy>/goai-s10-gate16-policy-v1-5/training/mujoco_s10/train_official_policy_residual.py}{train\_official\_policy\_residual.py}、\href{<workspace-legacy>/goai-s10-gate16-policy-v1-5/training/mujoco_s10/train.py}{train.py}。
\item
  \textbf{S6 --- 局部环境与奖励：}
  \href{<workspace-legacy>/goai-s10-gate16-policy-v1-5/training/mujoco_s10/official_policy_env.py}{official\_policy\_env.py}、\href{<workspace-legacy>/goai-s10-gate16-policy-v1-5/training/mujoco_s10/evaluate_official_track_skill.py}{evaluate\_official\_track\_skill.py}、\href{<workspace-legacy>/goai-s10-gate16-policy-v1-5/training/mujoco_s10/speed_objective.py}{speed\_objective.py}。
\item
  \textbf{S7 --- 发布模型选择：}
  \href{<workspace-legacy>/gate16_front_tuck_release_20260817/provenance/selection.json}{selection.json}、\href{<workspace-legacy>/goai-s10-gate16-policy-v1-5/docs/gate16_rear_push_result_20260817.md}{rear-push
  发布报告}、\href{<workspace-legacy>/goai-s10-gate16-policy-v1-5/docs/GATE16_HANDOFF_TRAINING_AND_REPO_20260817.md}{前期训练记录}。
\item
  \textbf{S8 --- 失败批次：}
  \href{<workspace-legacy>/resources/submission_mac_retest_20260820/20260820_215224_main3660b81/REPORT.md}{macOS
  submission retest}。
\item
  \textbf{S9 --- 依赖与模型来源：}
  \href{<workspace-legacy>/resources/submission_update_20260820_392s/submission_20260820_392s_main3660b81/03_documents/THIRD_PARTY.md}{THIRD\_PARTY.md}。
\item
  \textbf{S10 --- 感知与 SDK 部署实现：}
  \href{<workspace-legacy>/resources/submission_mac_retest_20260820/20260820_215224_main3660b81/source/goai26-s10-racing/src/s10_perception/s10_perception/lidar.py}{lidar.py}、\href{<workspace-legacy>/resources/submission_mac_retest_20260820/20260820_215224_main3660b81/source/goai26-s10-racing/src/s10_perception/s10_perception/heightmap.py}{heightmap.py}、\href{<workspace-legacy>/resources/submission_mac_retest_20260820/20260820_215224_main3660b81/source/goai26-s10-racing/src/s10_auto_nav/s10_auto_nav/local_planner.py}{local\_planner.py}、\href{<workspace-legacy>/resources/submission_mac_retest_20260820/20260820_215224_main3660b81/source/goai26-s10-racing/integration/gate16_policy_runner.hpp}{gate16\_policy\_runner.hpp}、\href{<workspace-legacy>/resources/submission_mac_retest_20260820/20260820_215224_main3660b81/source/goai26-s10-racing/integration/joint_command_owner.hpp}{joint\_command\_owner.hpp}。
\item
  \textbf{S11 --- 局部评测指标实现：}
  \href{<workspace-legacy>/goai-s10-gate16-policy-v1-5/training/mujoco_s10/front_retention.py}{front\_retention.py}、\href{<workspace-legacy>/goai-s10-gate16-policy-v1-5/training/mujoco_s10/speed_objective.py}{speed\_objective.py}。
\end{itemize}

\end{document}
